\section{The flow for two factors: the model case}\label{sec:flow}

This section treats a nonsplit extension of two stable bundles directly.
It introduces the three tools used in general in
Sections~\ref{sec:diagonal-flow}--\ref{sec:complete-spectrum}: comparison
of metrics along the flow, a corrected approximate solution with integrable
error, and matching of the flow with the extension orbit at late times.

Let $X$ be a compact connected Riemann surface, and retain $F,G$, the
nonzero harmonic extension form $\beta$, and the notation $q,V_q,c,C,z,R$
of Section~\ref{sec:two-factor}. Thus $E=E_1$ is the nonsplit extension
of the nonisomorphic stable equal-slope bundles $G$ by $F$. We use flow
time $\tau$ to distinguish it from the scaling parameter of the earlier
sections. Starting from any smooth Hermitian metric on $E$, consider
the normalized Hermitian Yang--Mills flow
\begin{equation}\label{eq:metric-flow}
 h^{-1}\partial_\tau h=-2(K_h-\kappa\id),\qquad
 \kappa=\frac{2\pi\mu(E)}{V}.
\end{equation}
The factor $2$ fixes the time normalization in the result below.

We use the classical convergence theorem for Yang--Mills flow on
a curve: the solution exists for all time, and in unitary gauges its
Chern connections converge to the Hermitian--Einstein connection on
the graded bundle $F\oplus G$
\cite{Daskalopoulos1992,Rade1992}. For smooth initial data the
convergence, modulo gauge, is smooth; the higher regularity follows
from parabolic regularity. In metric language there are holomorphic
isometries from $(E,h(\tau))$ to $(F\oplus G,h_0,\db^{(\tau)})$
with $\db^{(\tau)}\to\db_0$ smoothly. This analytic convergence
and the identification of the graded limit are inputs here.

The metric asymptotics of this flow have been studied much more
generally by Haiden--Katzarkov--Kontsevich--Pandit
\cite{HKKP2018}. Their construction uses harmonic extension forms,
a Green operator correction, and the order-preserving property of
the metric flow. We use these ingredients in the two-factor case,
together with a comparison begun at arbitrarily late times, to
determine the leading coefficient of the small Dolbeault eigenvalue.

\begin{theorem}\label{thm:flow-two-factor}
For every smooth initial metric in \eqref{eq:metric-flow}, exactly one
eigenvalue on trace-free endomorphisms tends to zero, and
\begin{equation}\label{eq:flow-rate}
 \lambda_1(h(\tau))=\frac{1+o(1)}{2\tau},\qquad
 \lambda_2(h(\tau))\ge\delta>0
\end{equation}
for sufficiently large $\tau$.

In the unitary identifications supplied by smooth convergence of the
connections, the first spectral projection converges in every
$L^2(h_0)\to C^k(h_0)$ norm to the projection onto $\CC q$.
Equivalently, the unit first eigensection, with its phase fixed by
positive inner product with $q$, tends smoothly to $q/\sqrt{V_q}$.
Its limit recovers the two summand projections by
\[
 P_F=\sqrt{V_q}\,u_\infty+\frac ar\id,\qquad P_G=\id-P_F.
\]
\end{theorem}

The proof separates metric comparison, a corrected approximate
solution, and a local description of the fixed extension orbit.

\subsection{Order comparison and the endomorphism spectrum}

If two metrics $h_1,h_2$ on a fixed holomorphic bundle satisfy
$e^{-\epsilon}h_1\le h_2\le e^\epsilon h_1$, then their induced
norms on endomorphisms satisfy
\[
 e^{-2\epsilon}|u|_{h_1}^2
 \le |u|_{h_2}^2\le e^{2\epsilon}|u|_{h_1}^2.
\]
Indeed, in an $h_1$-orthonormal frame diagonalizing $h_2$, the
coefficient of $|u_{ij}|^2$ is the ratio of its $i$th and $j$th
eigenvalues. The same bounds hold for endomorphism-valued $(0,1)$-forms,
since the base metric is fixed. The Dolbeault operator and its
trace-free domain are independent of the bundle metric, so min--max
gives
\begin{equation}\label{eq:metric-spectral-comparison}
 e^{-4\epsilon}\lambda_j(h_1)
 \le\lambda_j(h_2)\le e^{4\epsilon}\lambda_j(h_1).
\end{equation}
A positive scalar multiple of either metric leaves these eigenvalues
unchanged.

The order comparison for the metric flow has a useful version with
an error term. Suppose a smooth family $\widehat h(\tau)$ has residual
\begin{equation}\label{eq:flow-residual}
 \mathcal E_{\widehat h}
 :=\widehat h^{-1}\partial_\tau\widehat h
       +2(K_{\widehat h}-\kappa\id),\qquad
 \lVert\mathcal E_{\widehat h}\rVert_{L^\infty,\op,\widehat h}
 \le f(\tau).
\end{equation}
If $h$ is an exact solution and
$e^{-\epsilon}\widehat h(T)\le h(T)\le e^\epsilon\widehat h(T)$,
then
\begin{equation}\label{eq:flow-barriers}
 e^{-\epsilon-\int_T^\tau f(s)\,ds}\widehat h(\tau)
 \le h(\tau)\le
 e^{\epsilon+\int_T^\tau f(s)\,ds}\widehat h(\tau).
\end{equation}
To see this, multiply $\widehat h$ by
$e^{\pm\epsilon\pm\int_T^\tau f}$.
A spatially constant scalar does not change Chern curvature, and its
logarithmic derivative adds $\pm f\id$ to the residual. The two
metrics are therefore respectively a subsolution and a supersolution.
The maximum principle for Hermitian metrics gives the assertion;
this is the order comparison of
\cite[Section~4.2, Propositions~4.3 and~4.5]{HKKP2018}, with the
error measured in the normalization \eqref{eq:flow-residual}.

\subsection{An approximate solution with a small total error}

Use determinant-one diagonal scaling
\begin{equation}\label{eq:flow-diagonal}
 g_\rho=\rho^q
 =\diag(\rho^{b/r}\id_F,\rho^{-a/r}\id_G),\qquad
 h_\rho=g_\rho^*h_0.
\end{equation}
It still transports the extension entry to $\rho\beta$, and its
endomorphism spectrum is that of Theorem~\ref{thm:two-factor}.
The uncorrected curvature is $\kappa\id+\rho^2C$.
Its component $\rho^2z$ perpendicular to $q$ generally prevents
$h_\rho$ from following \eqref{eq:metric-flow} under a scalar
change of parameter.

Put
\begin{equation}\label{eq:flow-correction}
 v=Rz,\qquad k_\rho=\exp(-\tfrac12\rho^2v),\qquad
 \widehat g_\rho=k_\rho g_\rho,
 \qquad \widehat h_\rho=\widehat g_\rho^*h_0.
\end{equation}
The endomorphism $v$ is smooth, self-adjoint, trace-free, and block
diagonal, so it commutes with $q$. Block diagonality follows because
the split Laplacian preserves the diagonal blocks. Self-adjointness
follows from the K\"ahler identities: at a Hermitian--Einstein
connection the endomorphism Dolbeault Laplacian is half the connection
Laplacian and therefore preserves adjoints.

\begin{lemma}\label{lem:flow-approximation}
Let $0<\rho_0\le\rho_*$ for a fixed sufficiently small $\rho_*$,
and define, for $\tau\ge T$,
\begin{equation}\label{eq:flow-rho}
 \rho(\tau)=\bigl(\rho_0^{-2}+2c(\tau-T)\bigr)^{-1/2}.
\end{equation}
There is a constant $A$, independent of $T$ and $\rho_0$, such that
$\widehat h(\tau)=\widehat h_{\rho(\tau)}$ satisfies
\begin{equation}\label{eq:flow-error-bound}
 \lVert\mathcal E_{\widehat h}\rVert_{L^\infty,\op,\widehat h}
 \le A\rho(\tau)^3,
 \qquad
 \int_T^\infty A\rho(\tau)^3\,d\tau=\frac{A}{c}\rho_0.
\end{equation}
For this approximate solution,
\begin{equation}\label{eq:flow-approx-spectrum}
 \lambda_1(\widehat h_\rho)=c\rho^2(1+O(\rho^2)),\qquad
 \lim_{\tau\to\infty}\tau\lambda_1(\widehat h(\tau))=\tfrac12.
\end{equation}
\end{lemma}

\begin{proof}
Transport everything by $\widehat g_\rho$ to the fixed metric $h_0$.
At the split structure the derivative of mean curvature under the
metric change $h_0e^{\epsilon s}$ is $\Delta_0s$. For example, this
follows by differentiating the Chern curvature to obtain
$\sqrt{-1}\Lambda\db_0\partial_0s=\Delta_0s$.
Smooth Taylor expansion in $\rho$, with
$k_\rho=\id-\tfrac12\rho^2v+O_{C^k}(\rho^4)$, therefore gives
\begin{equation}\label{eq:flow-corrected-curvature}
 \widehat K_\rho
 :=\widehat g_\rho K_{\widehat h_\rho}\widehat g_\rho^{-1}
 =\kappa\id+\rho^2(C-\Delta_0v)+O_{C^k}(\rho^3)
 =\kappa\id+c\rho^2q+O_{C^k}(\rho^3).
\end{equation}
Terms involving both the change of metric, of order $\rho^2$,
and the extension, of order $\rho$, are included in the cubic
remainder. All estimates are uniform for $0<\rho\le\rho_*$.

Since $q$ and $v$ commute and $\dot\rho=-c\rho^3$, the logarithmic
derivative of the gauge is exactly
\[
 \dot{\widehat g}_\rho\widehat g_\rho^{-1}
 =-c\rho^2q+c\rho^4v.
\]
For a metric $g^*h_0$, its logarithmic time derivative, transported
by $g$, is $\dot g g^{-1}+(\dot g g^{-1})^*$.
Combining this with \eqref{eq:flow-corrected-curvature} cancels the
quadratic terms in \eqref{eq:flow-residual} and proves the cubic
bound. Integration using $\dot\rho=-c\rho^3$ gives the stated
total error.

Also $e^{-B\rho^2}h_\rho\le\widehat h_\rho\le e^{B\rho^2}h_\rho$
for some fixed $B$. Apply \eqref{eq:metric-spectral-comparison}
and Theorem~\ref{thm:two-factor} to get
$\lambda_1(\widehat h_\rho)=c\rho^2(1+O(\rho^2))$.
Equation~\eqref{eq:flow-rho} proves the limit.
\end{proof}

\subsection{Matching the actual flow at late times}

The following observation strengthens bounded comparison to comparison
with a factor tending to one. It uses the fact that the extension has
two nonisomorphic stable factors.

\begin{lemma}\label{lem:flow-matching}
For the solution $h(T)$ of \eqref{eq:metric-flow}, there are
$\rho_T>0$, $A_T>0$, and $\epsilon_T\ge0$, with
$\rho_T\to0$ and $\epsilon_T\to0$, such that
\begin{equation}\label{eq:flow-matching}
 e^{-\epsilon_T}A_T\widehat h_{\rho_T}
 \le h(T)\le e^{\epsilon_T}A_T\widehat h_{\rho_T}.
\end{equation}
\end{lemma}

\begin{proof}
Use the smooth convergence of the flow in unitary gauges, and write
$D_T=\db^{(T)}\to\db_0$ on $(F\oplus G,h_0)$, with each
$(F\oplus G,D_T)$ holomorphically isomorphic to $E$.
The spaces of holomorphic maps $F\to E$ and $E\to G$ both have
dimension one: apply $\Hom(F,\cdot)$ and $\Hom(\cdot,G)$ to the
extension and use stability and $\Hom(F,G)=0$. Their counterparts
at $F\oplus G$ are also one-dimensional, generated by the inclusion
and projection.

The Dolbeault Laplacians on these two Hom bundles converge smoothly.
The spectral projections onto their one-dimensional kernels therefore
converge in every smooth norm, by elliptic resolvent estimates.
Choose holomorphic maps $i_T:F\to(F\oplus G,D_T)$ and
$p_T:(F\oplus G,D_T)\to G$ converging respectively to inclusion
and projection. For large $T$, $i_T$ is an embedding and $p_T$ is
surjective. Moreover $p_Ti_T=0$, since $\Hom(F,G)=0$.

Set $j_T=p_T^*(p_Tp_T^*)^{-1}$ and
$\sigma_T(f,g)=i_Tf+j_Tg$. This smooth splitting tends to the
identity and has $p_T\sigma_T(f,g)=g$. Holomorphicity of $i_T,p_T$
implies
\[
 \sigma_T^{-1}D_T\sigma_T
 =\begin{pmatrix}\db_F&\alpha_T\\0&\db_G\end{pmatrix},
 \qquad \alpha_T\to0\quad\hbox{in }C^\infty.
\]
Hodge decomposition writes
$\alpha_T=H\alpha_T+\db v_T$, with $v_T\to0$ smoothly.
The unipotent change of splitting with upper-right entry $-v_T$
makes the extension entry $H\alpha_T$, while the resulting
isomorphism $S_T$ still tends to the identity.

The harmonic form $H\alpha_T$ is a nonzero scalar multiple
$z_T\beta$. Indeed, an isomorphism with $E$ identifies its inclusion
of $F$ and quotient map to $G$ up to nonzero scalars, because the
two Hom spaces above have dimension one. It therefore identifies
the extension classes up to their scalar ratio. Harmonic
representatives preserve this equality. Nonvanishing follows from
nonsplitting, and $z_T\to0$ follows from $\alpha_T\to0$.

Set $\rho_T=|z_T|$. A block-scalar unitary transformation removes
the phase of $z_T$. The pulled-back metric on
$(F\oplus G,\db_0+\rho_TN)$ tends to $h_0$ smoothly, since
$S_T\to\id$ and the phase transformation is unitary and parallel.
Its holomorphic identification with $(E,h(T))$ differs from
$g_{\rho_T}:E\to E_{\rho_T}$ by a scalar automorphism of $E$,
using simplicity from Theorem~\ref{thm:two-factor}. Consequently,
for some $A_T>0$,
\[
 h(T)=a_Tg_{\rho_T}^*\widetilde h_T,\qquad
 \widetilde h_T\longrightarrow h_0\quad\hbox{in }C^\infty.
\]
Thus $h(T)$ is within a multiplicative factor tending to one of
$a_Th_{\rho_T}$. The same is true of $\widehat h_{\rho_T}$
relative to $h_{\rho_T}$ by \eqref{eq:flow-correction}, which
proves \eqref{eq:flow-matching}.
\end{proof}

\begin{proof}[Proof of Theorem~\ref{thm:flow-two-factor}]
For a sufficiently large $T$, start the approximate solution of
Lemma~\ref{lem:flow-approximation} with $\rho_0=\rho_T$.
Multiply it by $A_T$, which changes neither its curvature nor its
endomorphism spectrum. Lemma~\ref{lem:flow-matching} supplies its
initial comparison with $h(T)$. Equations
\eqref{eq:flow-barriers} and \eqref{eq:flow-error-bound} give,
for every $\tau\ge T$,
\[
 e^{-b_T}A_T\widehat h(\tau)
 \le h(\tau)\le e^{b_T}A_T\widehat h(\tau),\qquad
 b_T=\epsilon_T+(A/c)\rho_T\longrightarrow0.
\]
It follows from \eqref{eq:metric-spectral-comparison} and
\eqref{eq:flow-approx-spectrum} that
\[
 \tfrac12e^{-4b_T}
 \le\liminf_{\tau\to\infty}\tau\lambda_1(h(\tau))
 \le\limsup_{\tau\to\infty}\tau\lambda_1(h(\tau))
 \le\tfrac12e^{4b_T}.
\]
Now let $T\to\infty$ to prove \eqref{eq:flow-rate} for the first
eigenvalue. A comparison begun at a single fixed time would give
only two-sided bounds of order $\tau^{-1}$, not the coefficient.

In the unitary gauges of smooth connection convergence, the
endomorphism Laplacians converge smoothly as operators with domain
$H^2$. Their limit is $\Delta_0$, whose kernel on trace-free
endomorphisms is $\CC q$. A contour enclosing only zero therefore
gives spectral projections of rank one for large $\tau$ and a
uniform gap above them. Resolvent convergence on the Sobolev scale
gives smooth convergence of the normalized eigensections and hence
of the finite-rank projections in $L^2\to C^k$ norm for every $k$.
This proves the remaining claims and the recovery formulas.
\end{proof}

The leading coefficient in \eqref{eq:flow-rate} is independent of
the ranks, extension norm, and initial metric under the normalization
\eqref{eq:metric-flow}. These data enter the intermediate constant
$c$, which cancels between the model spectrum $c\rho^2$ and the
flow law $\rho^2\sim(2c\tau)^{-1}$. The theorem does not assert
a next-order eigenvalue expansion or a quantitative rate of
eigensection convergence along the actual flow. In particular, the
fourth-order coefficient in \eqref{eq:two-expansion} cannot simply
be transferred by substituting $\rho=(2c\tau)^{-1/2}$ into the
uncorrected harmonic family. Sections~\ref{sec:diagonal-flow}
to~\ref{sec:complete-spectrum} extend the argument to every simple bundle on
a curve with pairwise nonisomorphic Jordan--H\"older factors, including
eigenvalues with distinct scales.
